\section{Concorrência e alternativas}

Uma das medidas efetuadas durante o planeamento da aplicação foi a análise de concorrência e o estudo das alternativas existentes. Desta forma foi possível adicionar mais funcionalidades à aplicação de forma a que esta fosse o mais completa possível.

\subsection{Concorrência} % (fold)
\label{sub:concorrencia}

De todas as aplicações concorrentes, as que mais se destacam são os Sistemas de Gestão de Aprendizagem (\emph{SGA}),também conhecidas por \emph{Learning Management Systems} - \emph{LMS}. Dentro dos \emph{SGA} imperam as seguintes aplicações:
\begin{itemize}
	\item \emph{Blackboard}
	\item \emph{Moodle}
	\item \emph{Studifi}
	\item \emph{Haiku LMS}
\end{itemize}

\subsubsection{Blackboard} % (fold)
\label{ssub:blackboard}

\begin{figure}[H]
        \centering
        \includegraphics[width=0.5\textwidth]{images/concorrencia/blackboard.jpg}
         \caption{\emph{Blackboard}}
         \label{fig: blackboard}
\end{figure}

A aplicação \href{http://www.blackboard.com}{\emph{Blackboard}} é composta por quatro tipos de utilizadores:
\begin{description}
	\item[Aluno] Um aluno pode submeter projetos, inscrever-se em grupos e consultar a informação disponibilizada pelos professores.
	\item[Professores] Um professor pode criar projetos e fazer avaliações qualitativas e quantitativas dos trabalhos enviados. Os projetos criados podem ter correção automática.
	\item[Técnicos] Os técnicos são responsáveis pela manutenção da aplicação.
	\item[Supervisores] Os supervisores fazem a gestão de cursos, disciplinas, docentes e alunos. Estes também podem gerar estatísticas do sistema.
\end{description}

Os projetos criados na \emph{Blackboard} não estão disponíveis para o público. Também não existe a noção de turnos, sendo estes substituídos pelos grupos. Os grupos são criados dentro da disciplina e um aluno pode estar inscrito em múltiplos grupos.

\subsubsection{Moodle} % (fold)
\label{ssub:moodle}
\begin{figure}[H]
        \centering
        \includegraphics[width=0.5\textwidth]{images/concorrencia/moodle.jpg}
         \caption{\emph{Moodle}}
         \label{fig: moodle}
\end{figure}
A aplicação \href{http://www.moodle.org}{\emph{Moodle}} é composta por três tipos de utilizadores:

\begin{description}
	\item[Alunos] Os alunos podem submeter projetos e consultar a informação disponibilizada pelos professores.
	\item[Professores] Os professores podem criar projetos e fazer uma avaliação quantitativa dos projetos enviados. Os projetos criados podem ter correção automática.
	\item[Administradors] Os Administradores são responsáveis pela manutenção da aplicação. Também são responsáveis por fazer a gestão dos utilizadores existentes.
\end{description}


Um dos problemas do \emph{Moodle} é a inexistência de grupos e de turnos. Quanto a projetos públicos apenas se estes forem inseridos manualmente pelos alunos nos seus portefólios.

\subsubsection{Studifi} % (fold)
\label{ssub:studifi}

\begin{figure}[H]
        \centering
        \includegraphics[width=0.5\textwidth]{images/concorrencia/studifi.png}
         \caption{\emph{Studifi}}
         \label{fig: studifi}
\end{figure}

A aplicação \href{https://studifi.com/}{\emph{Studifi}} é composta por três tipos de utilizadores:

\begin{description}
	\item[Alunos] Os alunos podem submeter projetos e consultar a informação disponibilizada pelos professores.
	\item[Professores] Os professores podem criar exercícios que são avaliados automaticamente e podem criar projetos.
	\item[Administradors] Os Administradores são responsáveis pela manutenção da aplicação. Também são responsáveis por fazer a gestão dos utilizadores existentes.
\end{description}

No \emph{Studifi} os exercícios não estão disponíveis para o público. Uma das características que se destaca nesta aplicação é a existência de controlo de versões nos ficheiros submetidos.


\subsubsection{Haiku LMS} % (fold)
\label{ssub:haiku_lms}

\begin{figure}[H]
        \centering
        \includegraphics[width=0.5\textwidth]{images/concorrencia/haiku.png}
         \caption{\emph{Haiku LMS}}
         \label{fig: studifi}
\end{figure}

A aplicação \href{http://www.haikulearning.com/}{\emph{Haiku LMS}} é composta por três tipos de utilizadores:

\begin{description}
	\item[Alunos] Os alunos podem submeter projetos e consultar a informação disponibilizada pelos professores.
	\item[Professores] Os professores podem criar projetos e exercícios e gerir pautas de avaliação.
	\item[Instituições de ensino] Responsáveis pela gestão de cursos e disciplinas da instituição.
\end{description}

\subsection{Alternativas} % (fold)
\label{sub:alternativas}

No que diz ao respeito às alternativas, estas são divididas em cinco situações:
\begin{itemize}
	\item Submissão de trabalhos

	\begin{description}
		\item[Submissão via \emph{mail}] O aluno envia um \emph{mail} com o trabalho para o \emph{mail} de um docente. Muitas vezes com um assunto especifico para que o docente possa filtrar o \emph{mail} e obter os vários trabalhos enviados.
		\item[Partilha de \emph{links}] O aluno faz \emph{upload} do seu trabalho para um servidor externo(\emph{github,dropbox,etc}) e partilha o \emph{link} com o docente.
		\item[Entrega de unidades de armazenamento] O aluno grava o seu trabalho numa unidade de armazenamento(\emph{Usb flash drive, cd/dvd, etc}) e entrega a respetiva unidade ao docente.
	\end{description}

	\item Gestão de pautas

	\begin{description}
		\item[Folhas de cálculo ou \emph{PDF}] É entre aos alunos um ficheiro com as notas do projeto. Muitas vezes a pauta entregue não refere a qualidade do trabalho.
		\item[Vitrinas] O docente imprime as notas finais e afixa-as numa vitrina do estabelecimento de ensino.
	\end{description}

	\item Correção automática

		\begin{description}
			\item[Script] Fica ao critério do docente, as funcionalidades desse emph{script}, isto é se o emph{script} criado verifica por exemplo extensões de ficheiros ou até mesmo se os ficheiros enviados correspondem a um determinado emph{output} para um certo \emph{input}.
		\end{description}
	\item Inscrições de grupos

	\begin{description}
		\item[\emph{Mail}] Um responsável do grupo envia um \emph{mail} ao docente com a constituição do grupo. Este \emph{mail} pode ficar sujeito as mesma regras que a submissão de trabalhos via \emph{mail}.
		\item[Entrega do trabalho] O grupo só fica registado e o docente só tem conhecimento desta após a entrega do trabalho.
		\item[Aviso com presença] Um elemento do grupo ou então todos os elementos notificam o docente em frente deste.
	\end{description}

	\item Gestão de grupos

	\begin{description}
		\item[Documento] O docente possui um documento com a constituição dos grupos e de outras informações relativas a esse grupo como por exemplo as notas dos trabalhos
	\end{description}

\end{itemize}




% subsection submiss_o_de_trabalhos (end)

% subsection alternativas (end)

\newpage
